\documentclass[11pt,a4paper]{article}

\usepackage[utf8]{inputenc}
\usepackage[T1]{fontenc}
\usepackage[spanish]{babel}
\usepackage{amsmath,amssymb}
\usepackage{geometry}
\geometry{margin=2.5cm}
\usepackage{xcolor}
\usepackage{array}
\usepackage{booktabs}
\usepackage{listings}
\usepackage{tikz}
\usepackage{enumitem}
\usepackage{hyperref}
\hypersetup{colorlinks=true, linkcolor=blue!50!black, urlcolor=blue!50!black}

% --- Estilo de código Python ---
\definecolor{codegray}{gray}{0.95}
\definecolor{keyword}{RGB}{0,0,180}
\definecolor{comment}{RGB}{0,128,0}
\definecolor{string}{RGB}{163,21,21}

\lstset{
  language=Python,
  backgroundcolor=\color{codegray},
  basicstyle=\ttfamily\footnotesize,
  keywordstyle=\color{keyword}\bfseries,
  commentstyle=\color{comment}\itshape,
  stringstyle=\color{string},
  showstringspaces=false,
  breaklines=true,
  frame=single,
  framesep=5pt,
  rulecolor=\color{gray!40},
  literate={í}{{\'i}}1 {á}{{\'a}}1 {é}{{\'e}}1 {ó}{{\'o}}1 {ú}{{\'u}}1 {ñ}{{\~n}}1 {Á}{{\'A}}1
}

\title{\textbf{Modelado matricial del embarque de un avión}}
\author{TP1 --- Aplicaciones Computacionales en Negocios}
\date{\today}

\begin{document}
\maketitle

\section{Idea general}

Simulamos el embarque de un avión de \textbf{25 filas} y \textbf{4 asientos por fila}
(100 pasajeros). Queremos medir cuánto tarda en llenarse según el \emph{orden de embarque}
(Random, Back-to-Front, WILMA, Steffen).

El modelo tiene dos piezas que conviene no confundir:

\begin{itemize}[leftmargin=1.4em]
  \item \textbf{Fuente de verdad:} los objetos \texttt{Pasajero}. Cada pasajero guarda
        sus atributos fijos y su \emph{historia} segundo a segundo.
  \item \textbf{Vista derivada:} la \texttt{matriz} del avión. Es una \emph{foto} del estado
        físico en el instante actual, recalculada en cada paso de tiempo.
\end{itemize}

La matriz \textbf{no} reemplaza a la clase \texttt{Pasajero}: la complementa. Sirve para
las preguntas \emph{espaciales} (¿quién bloquea a quién?, ¿cuántos hay a bordo?), que en
una lista de objetos serían incómodas de responder.

\section{Por qué no alcanza una sola matriz}

Una matriz de identificadores solo guarda \emph{posición}. Pero un pasajero carga más
información, y buena parte no es propia de una celda sino del pasajero:

\begin{center}
\begin{tabular}{l c}
\toprule
\textbf{Dato} & \textbf{¿Cabe en una matriz de \texttt{id}?} \\
\midrule
Posición actual & sí \\
Estado (camina / guarda / se sienta / \dots) & no (haría falta otra matriz) \\
Carry-on (0/1) & no --- es constante del pasajero \\
Velocidad, asiento destino & no --- ídem \\
Temporizador de la acción en curso & no \\
Historia por segundo & no --- la matriz es solo \emph{ahora} \\
\bottomrule
\end{tabular}
\end{center}

Conclusión: atributos e historia viven en \texttt{Pasajero}; la geometría del instante
vive en la \texttt{matriz} derivada. Al escribir la posición en un solo lugar, es
imposible que las dos representaciones se contradigan.

\section{La grilla}

\[
\texttt{matriz} \in \mathbb{Z}^{25 \times 5},
\qquad
\texttt{matriz[f, c]} =
\begin{cases}
0 & \text{celda vacía} \\
\texttt{id} > 0 & \text{pasajero ubicado en esa celda}
\end{cases}
\]

Como cada \texttt{id} es único y aparece en \textbf{una sola} celda mientras está a bordo,
la matriz dice al instante dónde está cada uno.

\subsection*{Ejes}

\textbf{Filas} (0--24 en NumPy): las 25 filas físicas. El avión numera de 1 a 25, así que
\[
\texttt{fila\_matriz} = \texttt{fila\_avion} - 1.
\]

\textbf{Columnas} (0--4): los cuatro asientos más el pasillo central.

\begin{center}
\begin{tabular}{c c c c c}
\toprule
col 0 & col 1 & col 2 & col 3 & col 4 \\
ventana izq. & pasillo izq. & \textbf{PASILLO} & pasillo der. & ventana der. \\
lado $-2$ & lado $-1$ & (se camina) & lado $1$ & lado $2$ \\
\bottomrule
\end{tabular}
\end{center}

La \textbf{columna 2} es el pasillo central: el único lugar donde un pasajero
\emph{camina}. Los sentados ocupan las columnas 0, 1, 3, 4.

\subsection*{Mapa lado $\rightarrow$ columna}

\begin{lstlisting}
def lado_a_col(lado):
    return {-2: 0, -1: 1, 1: 3, 2: 4}[lado]
\end{lstlisting}

\section{Ejemplo de foto}

En el instante $t$, el pasajero 7 ya está sentado en la fila 2, lado $-1$ (col 1),
y el pasajero 12 camina por el pasillo a la altura de la fila 2 (col 2):

\begin{center}
\begin{tabular}{r|c|c|c|c|c|}
\multicolumn{1}{r}{} & \multicolumn{1}{c}{\scriptsize c0} & \multicolumn{1}{c}{\scriptsize c1} & \multicolumn{1}{c}{\scriptsize c2} & \multicolumn{1}{c}{\scriptsize c3} & \multicolumn{1}{c}{\scriptsize c4} \\
\cline{2-6}
f0 & 0 & 0 & 0 & 0 & 0 \\ \cline{2-6}
f1 & 0 & \textbf{7} & \textbf{12} & 0 & 0 \\ \cline{2-6}
f2 & 0 & 0 & 0 & 0 & 0 \\ \cline{2-6}
\multicolumn{6}{c}{\scriptsize $\vdots$} \\
\end{tabular}
\end{center}

\section{La matriz como vista derivada}

La verdad es \texttt{pasajero.ubicacion[-1]} (la posición actual de cada uno). En cada
paso de tiempo se \textbf{borra} la matriz y se \textbf{reconstruye} leyendo a los
pasajeros:

\begin{lstlisting}
def actualizar_matriz(self):
    self.matriz[:] = 0                    # borro la foto anterior
    for p in self.pasajeros:
        fila, col = p.ubicacion[-1]       # donde esta ahora
        if fila >= 0:                     # -1 = fuera del avion, no se dibuja
            self.matriz[fila, col] = p.id
\end{lstlisting}

\section{Para qué sirve la matriz}

\begin{description}[leftmargin=1.4em, style=nextline]
  \item[Detectar bloqueos] El pasillo (col 2) admite un solo pasajero por celda.
        Para avanzar de la fila $f$ a la $f{+}1$:
        \begin{lstlisting}
if self.matriz[fila + 1, 2] == 0:   # pasillo libre adelante
    # avanzo
else:
    # bloqueado: espero (estado 2)
        \end{lstlisting}
  \item[Contar a bordo] Sin recorrer la lista de pasajeros:
        \begin{lstlisting}
np.count_nonzero(self.matriz > 0)
        \end{lstlisting}
  \item[Visualizar] Imprimir la matriz en cada paso da una animación del embarque.
\end{description}

\section{El cuello de botella}

Toda la demora del embarque nace en la columna 2. Un pasajero que guarda su carry-on en
la fila 10 \emph{tapa} a todos los que vienen atrás: ellos quedan en estado
\emph{bloqueado} hasta que el de adelante se siente y libere el pasillo. La representación
matricial hace ese choque explícito y trivial de chequear: basta mirar la celda de
adelante en la col 2.

\section{Cómo encaja con la simulación}

El orden de embarque (la salida de \texttt{orden\_random}, \texttt{orden\_steffen}, etc.)
es la \emph{cola}: define quién sube primero. El reloj avanza de a un segundo:

\begin{lstlisting}
t = 0
while not avion.esta_lleno():
    for p in pasajeros:            # en el orden de la cola
        p.actualizar_estado(avion, t)   # appendea 1 valor a ubicacion y estado
    avion.actualizar_matriz()      # recalculo la foto
    t += 1
\end{lstlisting}

Como cada \texttt{actualizar\_estado} agrega exactamente un valor por segundo, todas las
listas \texttt{ubicacion} y \texttt{estado} terminan con largo $t$: queda la película
completa del embarque, segundo a segundo, lista para medir tiempos o animar.

\end{document}
